% !TeX program = xelatex
% =====================================================================
%  Demo του θέματος beamerthemepdp.sty
% =====================================================================
\documentclass[t,aspectratio=169,9pt]{beamer}
\makeatletter\def\input@path{{../../theme/}}\makeatother   % το κοινό θέμα βρίσκεται στο Presentations/theme/
\usetheme{pdp}

\title{Δομή Επανάληψης}
% \subtitle{Demo του νέου θέματος}
\author{Αναστάσης Κολιόπουλος}
% \institute{Πανελλήνιος Διαγωνισμός Πληροφορικής}
\date{}

\begin{document}

\pdptitleframe
\pdptocframe{Περιεχόμενα}

% =====================================================================
\section{Ο βρόχος \texttt{while}}

\begin{frame}{Ο βρόχος \texttt{while}}
  \field{Είσοδος:}
  \pause
  \begin{itemize}
    \item Κατευθυνόμενο γράφημα $G = (V, E)$ με $w(e) \geq 0$ για κάθε $e \in E$, όπου $n = |V|$ και $m = |E|$.
    \pause
    \item Αφετηριακή κορυφή $s \in V$.
  \end{itemize}
  \pause
  \field{Έξοδος:}
  \pause
  \begin{itemize}
    \item Η απόσταση $d(v)$ από την $s$ για κάθε κορυφή $v \in V$.
  \end{itemize}
  \pause
  \field{Ιδέα:}
  \pause
  Σε κάθε βήμα «οριστικοποιούμε» την κορυφή με τη μικρότερη προσωρινή απόσταση $D[v]$ και χαλαρώνουμε τις ακμές της:
  \[
    D[u] \gets \min\bigl\{D[u],\; D[v] + w(v, u)\bigr\}, \qquad \forall u \in N^+(v).
  \]
\end{frame}



\begin{frame}{Ο βρόχος \texttt{while}}
  \begin{block}{Παρατήρηση.}
    \pause
    Όταν μία κορυφή $v$ εξάγεται από την ουρά προτεραιότητας, ισχύει $D[v] = d(v)$.
  \end{block}
  \pause
  \begin{block}{\emph{Απόδειξη.}}
    Έστω $v$ η πρώτη κορυφή που εξάγεται με $D[v] > d(v)$ και $P$ ένα συντομότερο $s - v$ μονοπάτι.
    \pause
    Έστω $(x, y)$ η πρώτη ακμή του $P$ με $x$ ήδη εξαγμένη και $y$ όχι. Τότε
    \[
      D[y] \leq D[x] + w(x, y) = d(y) \leq d(v) < D[v],
    \]
    \pause
    αφού τα βάρη είναι μη αρνητικά. Άρα η $y$ θα είχε εξαχθεί πριν από τη $v$ (άτοπο).
    \hfill$\blacksquare$
  \end{block}
\end{frame}

\begin{frame}{Ο αλγόριθμος του Dijkstra}{Ψευδοκώδικας}
  \begin{columns}[T]
    \begin{column}{0.52\textwidth}
      \begin{pseudo}
        \Procname{$\proc{Dijkstra}(G, w, s)$}
        \li $\proc{Initialize-Single-Source}(G, s)$
        \li $S \gets \emptyset$
        \li $Q \gets \attrib{G}{V}$
        \li \While $Q \neq \emptyset$
        \li   \Do
                $u \gets \proc{Extract-Min}(Q)$
        \li     $S \gets S \cup \{u\}$
        \li     \For κάθε κορυφή $v \in \attrib{G}{Adj}[u]$
        \li       \Do
                    $\proc{Relax}(u, v, w)$ \Comment{χαλάρωση}
                  \End
              \End
        \li   \For $i == 1; \; i \leq n; \; i++$
        \li     \Do
        \li ssss
              \End
      \end{pseudo}
    \end{column}
    \begin{column}{0.44\textwidth}
      \pause
      \begin{itemize}
        \item Η $Q$ είναι ουρά προτεραιότητας με κλειδί το $\attrib{v}{d}$.
        \pause
        \item Κάθε κορυφή εξάγεται μία φορά (γραμμή 5).
        \pause
        \item Κάθε ακμή χαλαρώνεται μία φορά (γραμμή 8).
      \end{itemize}
    \end{column}
  \end{columns}
\end{frame}

\begin{frame}[fragile]{Ο αλγόριθμος του Dijkstra}{Υλοποίηση σε C++}
  \begin{columns}[T]
    \begin{column}{0.62\textwidth}
\begin{cpp}
using ll = long long; using P = pair<ll, int>;
vector<ll> dijkstra(int s) {     // adj: γείτονες
  vector<ll> D(n, INF);
  priority_queue<P, vector<P>, greater<>> pq;
  D[s] = 0; pq.push({0, s});
  while (!pq.empty()) {
    auto [d, v] = pq.top(); pq.pop();
    if (d > D[v]) continue;      // παλιά εγγραφή
    for (auto [u, w] : adj[v])
      if (D[v] + w < D[u]) {
        D[u] = D[v] + w;          // χαλάρωση
        pq.push({D[u], u});
      }
  }
  return D;
}
\end{cpp}
    \end{column}
    \begin{column}{0.34\textwidth}
      \pause
      \begin{itemize}
        \item Η \texttt{priority\_queue} με \texttt{greater<>} δίνει πάντα την κορυφή με το μικρότερο $D[v]$.
        \pause
        \item Αντί για \emph{decrease-key}, βάζουμε νέα εγγραφή και αγνοούμε τις παλιές με το \texttt{continue}.
        \pause
        \item Συνολικά $O\bigl((n + m) \lg n\bigr)$.
      \end{itemize}
    \end{column}
  \end{columns}
\end{frame}

\begin{frame}{Ο αλγόριθμος του Dijkstra}
  \begin{block}{Χρονική πολυπλοκότητα.}
    \begin{itemize}
      \item $O(n \lg n + m \lg n)$ - \emph{Δυαδικός σωρός}.
      \item $O(n \lg n + m)$ - \emph{Σωρός Fibonacci}.
    \end{itemize}
  \end{block}
  \pause
  \begin{alertblock}{Προσοχή.}
    \pause % πρώτα μόνο το κουτί με τον τίτλο, το κείμενο στο επόμενο βήμα
    Με αρνητικά βάρη ο αλγόριθμος δεν είναι πια σωστός· χρειαζόμαστε τον Bellman--Ford.
  \end{alertblock}
  \pause
  \vspace{0.3cm}
  \begin{center}
    \pdpcodelink{https://github.com/anastasiskolio13/ShortestPaths/blob/master/Traveller_Dijkstra.cpp}{Ενδεικτική κωδικοποίηση}
  \end{center}
\end{frame}

% =====================================================================
\section{Ο βρόχος \texttt{do while}}

\begin{frame}{Ο βρόχος \texttt{do while}}
  Ορίζουμε ως $dp[i][v]$ το μήκος του συντομότερου $s - v$ μονοπατιού με το πολύ $i$ ακμές.
  \pause
  \[
    dp[i][v] =
    \begin{cases}
      0, & i = 0 \text{ και } v = s, \\
      \infty, & i = 0 \text{ και } v \neq s, \\
      \min\Bigl\{dp[i - 1][v],\; \displaystyle\min_{u \in N^-(v)}\bigl\{dp[i - 1][u] + w(u, v)\bigr\}\Bigr\}, & i \geq 1.
    \end{cases}
  \]
  \pause
  \begin{exampleblock}{Παρατήρηση.}
    Αν δεν υπάρχει αρνητικός κύκλος, κάθε συντομότερο μονοπάτι έχει το πολύ $n - 1$ ακμές,
    \pause % το κουτί εμφανίζεται μαζί με την πρώτη φράση
    άρα η απάντηση είναι το $dp[n - 1][v]$.
  \end{exampleblock}
  \pause
  Αρκεί να κρατάμε μόνο τις δύο τελευταίες γραμμές του πίνακα: $dp[i \bmod 2][v]$.
\end{frame}

\begin{frame}[fragile]{Ο βρόχος \texttt{do while}}{Από το αρχείο \texttt{BellmanFord.cpp}}
  \cppfile[firstline=35, lastline=46]{BellmanFord.cpp}
  \pause
  \begin{block}{Χρονική πολυπλοκότητα.}
    $O(n m)$ χρόνος και $O(n)$ μνήμη.
  \end{block}
\end{frame}

\section{Ο βρόχος \texttt{for}}

\begin{frame}{Βιβλιογραφία}
  \begin{thebibliography}{99}
    \bibitem{dmagos} Δημήτριος Μάγος. \textit{Στοιχεία Ανάλυσης Αλγορίθμων}. Εκδόσεις Τσότρας, 2024.
    \bibitem{cormen2022} Cormen, Leiserson, Rivest, Stein. \textit{Introduction to Algorithms}. 4η Έκδοση, MIT Press, 2022.
  \end{thebibliography}
\end{frame}

\pdpclosingframe

\end{document}
